\section{A predecessor from equations and equation induction alone}
\label{app:pure-equational-predecessor}

This appendix supplies the arithmetic fact needed by the restricted-witness
theorem: successor injectivity follows from a purely equational fragment,
without case introduction for $D$.  The construction of the predecessor is inherited
from the repository's scalar bootstrap.  The contribution of this appendix
is its purely equational organization, in particular the replacement of
case splits made under an induction hypothesis by unconditional finite
identities proved in advance.  The argument below uses no numerical
coverage axiom, no logical disjunction, and no successor injectivity.
After equation induction is made available through the interface below,
the arithmetic argument itself has no $D$ formulas.

\subsection{Resources and a finite-value lemma}

The available terms are variables, $0$, $S$, and the iterator
$I(n,a;u.h)$, whose only bound variable is $u$ in $h$.  We use
\begin{align}
 I(0,a;u.h)&=a, \label{eq:pe-zero}\\
 I(Sn,a;u.h)&=h[I(n,a;u.h)/u], \label{eq:pe-succ}
\end{align}
ordinary equality reasoning, and induction on a single equation.  In the
native calculus the latter is obtained from QI by progress and QE:
\[
 \frac{\Gamma\vdash E(0)\qquad
       \Gamma,E(z)\vdash E(Sz)}{\Gamma\vdash E(t)},
 \qquad z\notin\operatorname{FV}(\Gamma),
\]
where $E$ is a literal equation and all its other parameters stay fixed.
In a combined-introduction presentation, QI by progress is obtained by
introducing an equation matrix with unused witness at $0$ and then opening
that binder.  The equational argument itself needs only the displayed
equation-induction interface.

\paragraph{Instantiation in the equational fragment.}
Capture-avoiding substitution of arbitrary terms for numerical parameters
is admissible in this equational fragment.  This follows by induction on
an equational derivation: computation and transport commute with
substitution, and an induction eigenvariable is first renamed fresh for
the terms being substituted.  In the combined-rule implementation of
equation induction, the witness remains the literal $0$ and the matrix
remains an equation.  Thus this argument never substitutes into a case
introduction at a variable witness.  Completed equational identities may
therefore be instantiated at arbitrary terms without presupposing a
substitution theorem for the full restricted calculus.

Ignoring the induction hypothesis gives \emph{equational cases}: proofs of
$E(0)$ and $E(Sz)$ give $E(t)$ for every term $t$.  This is not a case rule
for arbitrary formulas.  It is a derived rule only for equations.

Write $1=S0$, with larger numerals defined similarly, and put
\[
 C(n,a,b)=I(n,a;u.b),\qquad
 Z(n)=C(n,1,0),\qquad
 N(n)=I(n,0;u.Z(u)).
\]
The update variable is absent from the body of $C$, so computation gives
\begin{gather}
 C(0,a,b)=a,\qquad C(Sn,a,b)=b,\label{eq:pe-selector}\\
 Z(0)=1,\qquad Z(Sn)=0,\label{eq:pe-zero-test}\\
 N(0)=0,\qquad N(Sn)=Z(N(n)).\label{eq:pe-parity-rec}
\end{gather}

\paragraph{Finite-value transfer.}
Let $F(v)$ be an equation, and suppose that $F(0)$ and $F(1)$ have been
proved over a fixed side context $\Gamma$.  Then $F(Z(t))$ and $F(N(t))$
are derivable over $\Gamma$ for every term $t$.  The placeholder $v$ is
fresh for $\Gamma$; all other parameters of $F$ remain fixed.

For the first conclusion, use equational cases on a fresh variable $x$
with equation $F(Z(x))$.  Its zero case is $F(1)$ and its successor case is
$F(0)$, by \eqref{eq:pe-zero-test}.  For the second, use equational cases
on $x$ with equation $F(N(x))$.  Its zero case is $F(0)$; its successor
case is $F(Z(N(x)))$, already provided by the first conclusion.
Finally instantiate the completed equations at $x=t$.

In subsequent uses of this lemma, the two equations $F(0)$ and $F(1)$
are unconditional computations.  We never infer an equation by splitting
on $N(t)=0$ or $N(t)=1$ while a side assumption depends on that value.

\subsection{Addition and parity identities}

Define
\[
 a+b=I(b,a;u.Su),\qquad 2a=a+a,\qquad 3a=a+a+a.
\]
Thus $a+0=a$ and $a+Sb=S(a+b)$.  Equation induction on $b$ gives
$0+b=b$ and $Sa+b=S(a+b)$.  Another induction on $b$ proves
commutativity: the successor step is
\[
 a+Sb=S(a+b)=S(b+a)=Sb+a.
\]
Induction on the last summand proves associativity.  We shall reorder
finite sums using these completed identities.  In particular,
$a+1=Sa$, $2(Sa)=2a+2$, and $3(Sa)=3a+3$.

The instances at $v=0,1$ of
\[
 v+Z(v)=1,\qquad Z(Z(v))=v
\]
compute directly.  Finite-value transfer therefore gives
\begin{equation}
 N(n)+Z(N(n))=1,\qquad Z(Z(N(n)))=N(n).
 \label{eq:pe-parity-bool}
\end{equation}
Two successor computations in \eqref{eq:pe-parity-rec} now give
\begin{equation}
 N(n+2)=N(n).\label{eq:pe-period-two}
\end{equation}
An equation induction on $n$ proves $N(2n)=0$: the successor step reduces
$N(2n+2)$ to $N(2n)$.  Consequently $N(2n+1)=Z(0)=1$.
An independent equation induction proves $N(3n)=N(n)$: after increasing
$n$, the left argument increases by three, so period two reduces the
new left value to $Z(N(3n))$, and the induction equation changes this to
$Z(N(n))=N(Sn)$.  We have proved
\begin{equation}
 N(2n)=0,\qquad N(2n+1)=1,\qquad N(3n)=N(n).
 \label{eq:pe-parity-scaling}
\end{equation}
No assertion that every number is even or odd has been used.

\subsection{A three-state update and a ramp}

Put
\[
 j(u)=Z(u)+2N(u),\qquad
 M(n)=I(n,0;u.j(u)),\qquad
 G(n)=I(n,0;u.u+M(u)+1).
\]
Finite computations give
\begin{equation}
 j(0)=1,\qquad j(1)=2,\qquad j(2)=0.
 \label{eq:pe-three-cycle}
\end{equation}
We need these values and $M(Sn)=j(M(n))$; we do not assume a remainder
theorem for $M$.

For arbitrary parameter $g$, the following equations hold whenever
$b=N(v)$:
\begin{align}
 M(g+b+1)&=C\bigl(b,j(M(g)),j(j(M(g)))\bigr),
    \label{eq:pe-ramp-finite-a}\\
 C\bigl(b,j(b),j(j(b))\bigr)&=Z(b).
    \label{eq:pe-ramp-finite-b}
\end{align}
To prove \eqref{eq:pe-ramp-finite-a}, apply finite-value transfer to the
equation with placeholder $b$.  At $b=0$ it is
$M(g+1)=j(M(g))$, and at $b=1$ it is
$M(g+2)=j(j(M(g)))$, both computation identities.  For
\eqref{eq:pe-ramp-finite-b}, the zero instance is
$j(0)=Z(0)=1$, and the one instance is
$j(j(1))=j(2)=Z(1)=0$.  These are again computations.  These two
identities are proved \emph{before} the next induction.

We claim
\begin{equation}
 M(G(n))=N(n).\label{eq:pe-ramp-invariant}
\end{equation}
The zero instance is immediate.  In the successor step retain the single
induction equation $M(G(n))=N(n)$ and calculate
\begin{align*}
 M(G(Sn))
 &=M\bigl(G(n)+M(G(n))+1\bigr)\\
 &=M\bigl(G(n)+N(n)+1\bigr)\\
 &=C\bigl(N(n),j(M(G(n))),j(j(M(G(n))))\bigr)\\
 &=C\bigl(N(n),j(N(n)),j(j(N(n)))\bigr)\\
 &=Z(N(n))=N(Sn).
\end{align*}
The second and fourth lines use the induction equation; the third and
fifth use the two unconditional identities just proved.  In particular,
there is no case split under the induction hypothesis.

It follows that
\begin{equation}
 G(Sn)=G(n)+N(n)+1,\qquad G(n+2)=G(n)+3.
 \label{eq:pe-ramp-step}
\end{equation}
The second equation follows from the first applied twice and
$N(n)+Z(N(n))=1$.  Equation induction gives
\begin{equation}
 G(2n)=3n.\label{eq:pe-ramp-double}
\end{equation}
For the successor step, use $G(2n+2)=G(2n)+3$.

A separate equation induction proves
\begin{equation}
 G(3n)=3G(n)+N(n).\label{eq:pe-ramp-triple}
\end{equation}
Indeed, \eqref{eq:pe-ramp-step} gives
$G(x+3)=G(x)+N(x)+4$.  Assuming the induction equation and using
$N(3n)=N(n)$, the new left side is
\[
 G(3Sn)=G(3n)+N(3n)+4=3G(n)+2N(n)+4.
\]
The new right side is
\[
 3G(Sn)+N(Sn)=3G(n)+3N(n)+3+Z(N(n)).
\]
The two expressions are equal by $N(n)+Z(N(n))=1$ and the additive
identities.  The base is $0=0$.  Again, no conditional case distinction
occurs inside the induction.

\subsection{The factor map}

Define
\[
 H(u)=C(N(u),G(u),u).
\]
Then
\begin{equation}
 H(2u)=3u,\qquad H(1)=1,\qquad H(3u)=3H(u).
 \label{eq:pe-factor}
\end{equation}
The first two equations follow immediately from
\eqref{eq:pe-parity-scaling}, \eqref{eq:pe-ramp-double}, and the two
selector computations.  For the third, first prove by finite-value
transfer the unconditional equation
\begin{equation}
 C(b,3a+b,3c)=3C(b,a,c),\qquad b=N(v).
 \label{eq:pe-factor-finite}
\end{equation}
At $b=0$ both sides are $3a$, and at $b=1$ both sides are $3c$.
Now substitute $v=u$, $a=G(u)$, and $c=u$ into this completed identity:
\begin{align*}
 H(3u)
 &=C\bigl(N(u),3G(u)+N(u),3u\bigr)\\
 &=3C(N(u),G(u),u)=3H(u).
\end{align*}
This argument does not represent an arbitrary number as $2v$ or $2v+1$.

\subsection{A comparator}

For a fixed update term $h$, whose parameters do not include the variable
being inducted on, equation induction proves the rotation identity
\begin{equation}
 I(Sy,a;h)=I(y,h(a);h).
 \label{eq:pe-rotation}
\end{equation}
At $y=0$ both sides are $h(a)$.  In the step, the two successor
computations apply the same fixed $h$ to the induction equation.
Here $I(y,a;h)$ abbreviates $I(y,a;u.h(u))$; it is a term abbreviation,
not quantification over functions.

Using $H(3u)=3H(u)$, another equation induction on $y$, with $a$ fixed,
gives
\begin{equation}
 I(y,3a;H)=3I(y,a;H).
 \label{eq:pe-orbit-scaling}
\end{equation}
The base is $3a=3a$; the step rewrites
$H(I(y,3a;H))$ to $H(3I(y,a;H))=3H(I(y,a;H))$.
Similarly, $H(1)=1$ gives
\begin{equation}
 I(y,1;H)=1.\label{eq:pe-orbit-one}
\end{equation}

Set
\[
 A(x)=I(x,1;u.2u),\qquad
 L(x,y)=N\bigl(I(y,A(x);H)\bigr).
\]
We obtain the three comparison equations
\begin{equation}
 L(0,y)=1,\qquad L(Sx,0)=0,\qquad L(Sx,Sy)=L(x,y).
 \label{eq:pe-comparator}
\end{equation}
The first uses $A(0)=1$, \eqref{eq:pe-orbit-one}, and $N(1)=1$.
The second uses $A(Sx)=2A(x)$ and $N(2u)=0$.  For the third,
\begin{align*}
 L(Sx,Sy)
 &=N\bigl(I(Sy,2A(x);H)\bigr)\\
 &=N\bigl(I(y,H(2A(x));H)\bigr)\\
 &=N\bigl(I(y,3A(x);H)\bigr)\\
 &=N\bigl(3I(y,A(x);H)\bigr)\\
 &=N\bigl(I(y,A(x);H)\bigr)=L(x,y).
\end{align*}
The successive steps use rotation, $H(2u)=3u$, orbit scaling, and
$N(3u)=N(u)$.  Thus \eqref{eq:pe-comparator} holds internally at all
arguments.  The intended interpretation of $L(x,y)$ as the indicator of
$x\leq y$ is not a premise.

\subsection{Saturation and predecessor}

Define the following open terms, with $m$ a fixed parameter:
\[
 B_m(u)=C(L(Su,m),u,Su),\qquad
 W(t,m)=I(t,0;u.B_m(u)).
\]
Equational cases on a fresh variable $e$ prove
\begin{equation}
 C(e,Sa,Sb)=S C(e,a,b).
 \label{eq:pe-selector-succ}
\end{equation}
At $e=0$ both sides are $Sa$, and at $e=Sz$ both are $Sb$.
Using \eqref{eq:pe-comparator} and this selector equation gives
\begin{align}
 B_{Sm}(Su)
 &=C(L(SSu,Sm),Su,SSu)\notag\\
 &=C(L(Su,m),Su,SSu)=S B_m(u),
 \label{eq:pe-saturation-conjugacy}\\
 B_{Sm}(0)
 &=C(L(S0,Sm),0,S0)\notag\\
 &=C(L(0,m),0,S0)=C(1,0,S0)=S0.
 \label{eq:pe-saturation-base}
\end{align}

Equation induction on $t$, with $m$ fixed, now proves
\begin{equation}
 W(St,Sm)=S W(t,m).
 \label{eq:pe-saturation-shift}
\end{equation}
For $t=0$, the left side is $B_{Sm}(0)=S0$, as required.  The step is
\begin{align*}
 W(SSt,Sm)
 &=B_{Sm}(W(St,Sm))\\
 &=B_{Sm}(S W(t,m))\\
 &=S B_m(W(t,m))=S W(St,m).
\end{align*}
Only the second line uses the active induction equation.

Two more equation inductions give
\begin{equation}
 L(Sn,n)=0,\qquad W(n,n)=n.
 \label{eq:pe-saturation-diagonals}
\end{equation}
For the first, the base is $L(S0,0)=0$ and the step reduces
$L(SSn,Sn)$ to $L(Sn,n)$.  For the second, the base is $W(0,0)=0$,
and the step uses the completed identity
\eqref{eq:pe-saturation-shift} at $t=m=n$:
$W(Sn,Sn)=S W(n,n)=Sn$.  This substitution is into an already proved
identity, and does not change a parameter of an active induction
hypothesis.  Consequently
\begin{equation}
 W(Sn,n)=B_n(n)=C(L(Sn,n),n,Sn)=n.
 \label{eq:pe-saturation-extra}
\end{equation}

Finally define the finite scalar term
\begin{equation}
 P(m)=I\bigl(m,0;u.C(L(SSu,m),u,Su)\bigr).
 \label{eq:pe-predecessor-term}
\end{equation}
Its zero equation is a computation.  For its successor equation put
\[
 V(t,m)=I\bigl(t,0;u.C(L(SSu,Sm),u,Su)\bigr).
\]
The third equation in \eqref{eq:pe-comparator} gives the pointwise
identity
\[
 C(L(SSu,Sm),u,Su)=C(L(Su,m),u,Su)=B_m(u).
\]
Equation induction on $t$, with $m$ fixed, proves $V(t,m)=W(t,m)$.
The base is $0=0$.  In the step, first apply the pointwise update
identity at $u=V(t,m)$ and then the induction equation, obtaining
$B_m(V(t,m))=B_m(W(t,m))$.  Thus no primitive extensionality rule for
iterator bodies is being assumed.  From
\eqref{eq:pe-saturation-extra} we conclude
\begin{equation}
 \boxed{\quad P(0)=0,\qquad
 P(Sm)=V(Sm,m)=W(Sm,m)=m.\quad}
 \label{eq:pe-predecessor-equations}
\end{equation}
Applying equality congruence for $P$ to $Sx=Sy$ and using
\eqref{eq:pe-predecessor-equations} gives $x=y$.

Every abbreviation in this appendix expands to a finite term over
$0,S,I$.  Every induction invariant is one literal equation.  Equality
reflection, pairing, general primitive-recursion compilation, numerical
coverage, and logical case analysis were not used to obtain predecessor
or successor injectivity.

\subsection{Provenance and verification status}

The arithmetic construction, including the factor map and the saturated
iteration, comes from
\texttt{Research/certified\_zero/reference/scalar\_predecessor.tex} in the
user's ConciseIPRA project, read at commit
\texttt{9f9545f5ffd9e0392d4a36639503aca12383450a}.
The repository URL is
\begin{center}\small
\href{https://github.com/VladimirReshetnikov/ProveIt/tree/9f9545f5ffd9e0392d4a36639503aca12383450a/Logic/PrimitiveRecursiveArithmetic/ConciseIPRA}{ConciseIPRA at the pinned commit}
\end{center}
The existing source explicitly allowed reconstructed numerical cases.
The present appendix strengthens its dependency statement by giving
equation-only proofs of the needed finite identities, particularly
\eqref{eq:pe-ramp-finite-a}, \eqref{eq:pe-ramp-finite-b}, and
\eqref{eq:pe-factor-finite}.  These identities remove the circularity that
would arise if a proof of restricted logical introduction already assumed
unrestricted logical case distinction.

This appendix is a written internal derivation.  It does not claim a newly
compiled Lean formalization of the equational dependency refinement or of
the restricted-witness theorem.  The repository's pre-existing Lean
formalization should not be presented as machine verification of this
new dependency claim without the corresponding new formal proof.
